\begin{textsample}{NEG Dim 1 – vem\_ed\_33 – Score 4.00 – vem\_ed\_33/t179.txt}  \label{ex:f1_neg_108}
Modelos \textbf{preditivos} para intrusões de computadores O PCI realizado no segundo semestre de 2025 entre estudantes de Big Data para Negócios da Fatec Ipiranga e de Redes e Sistemas Informáticos do ISP Gaya criou modelos \textbf{preditivos} para \textbf{detectar} intrusões em redes de computadores. “Criação de modelo \textbf{preditivo} através de sistema de detecção de intrusões em redes de computadores” foi o PCI orientado pelos professores Carlos Menezes e José Arnaud (ISP Gaya).
O contato com a instituição portuguesa foi realizado por Patrícia Patrício, da AIES / Depes / CGESG, que acompanhou desde as reuniões iniciais com os professores parceiros até o design e acompanhamento do projeto.
Os alunos do ISP Gaya elaboraram \textbf{relatórios} de tráfego normal e de intrusão de dados em um sistema controlado (Snort) e, a partir da análise \textbf{exploratória} dos dados, os estudantes da Fatec Ipiranga criaram modelos \textbf{preditivos}, com acurácia próxima de 100 \%. “Foram oito semanas de descobertas, de superação e de aprendizagem prática.
Mas mais do que um projeto, foi uma experiência humana, de trabalho em equipe, partilha de conhecimento e crescimento conjunto”, comentou José Arnaud, em postagem no LinkedIn. “O professor Arnaud é um parceiro \textbf{fantástico}; ele é organizado, acessível, gosta de ouvir e opinar, o que facilita muito a colaboração”, \textbf{elogia} o professor Menezes. “Para os estudantes, a sensação de interagir numa equipe real, com as dificuldades inerentes a atuar em equipes multicêntricas, foi o grande benefício do projeto.
Além disso, houve uma troca de informações: nossos alunos aprenderam sobre redes e os de Portugal sobre análise de dados”, conclui.

% matched lemmas: detectar, elogiar, exploratório, fantástico, preditivo, relatório
\end{textsample}
